\section{Literature review}\label{sec:lit}
We review four streams of research: train timetabling with stop planning, the train platforming problem, the integration of timetabling with station track and route assignment, and large neighborhood search and MILP-based matheuristics in railway scheduling. We then summarize the research gap.

\subsection{Train timetabling and stop planning}
The train timetabling problem determines arrival and departure times and train orders under shared infrastructure capacity; \citet{Cacchiani2012} survey its nominal and robust variants. Decomposition by train is a classical approach: \citet{Brannlund1998} relax the capacity constraints and solve train-level dynamic programs, and \citet{Caprara2002} combine a directed multigraph formulation with Lagrangian relaxation and a heuristic. Event-based formulations represent the timetable as a set of events linked by activities with minimum and maximum durations; the periodic event scheduling problem of \citet{Serafini1989} is the standard periodic version, and \citet{Liebchen2007} show how many practical requirements can be modelled on such networks. Time--space network formulations are a common alternative, for example for cyclic timetabling \citep{ZhangY2019b} and for integrated line planning and timetabling \citep{ZhangY2022}.

Flexible stopping patterns add service decisions to the timing problem. \citet{Yang2016} embed train stop planning, including station service frequency constraints, into a MILP train scheduling model for a high-speed corridor. \citet{Yue2016} optimize stopping patterns and schedules with a column-generation-based heuristic, and \citet{Cacchiani2020} propose robust models for integrated stop planning and timetabling under demand uncertainty. \citet{Dong2020} jointly optimize stop plans, timetables, and the number of trains for commuter railways under train and station service frequency constraints, using an extended ALNS. \citet{Yao2023} integrate stop planning and passenger routing into network periodic timetabling on a periodic time--space network, and \citet{Tian2024} additionally select train types through a variable-splitting Lagrangian decomposition. Recent studies in this journal address demand-driven timetabling with flexible periodicity solved by ALNS \citep{Yin2024} and timetabling with coupling and decoupling operations solved by a matheuristic \citep{Liu2024}. These studies treat stations at the macroscopic level, i.e., through station headways or capacities, without assigning tracks and routes.

\subsection{Train platforming problem}
The train platforming problem assigns each train an arrival--departure track (platform) together with its arrival and departure paths through the station throat area \citep{Caprara2011}. Track availability and route compatibility must both be checked, since train movements on different tracks may still use conflicting routes. \citet{Zwaneveld1996} formulate routing through stations as a node-packing problem and solve it by branch-and-cut, \citet{Zwaneveld2001} strengthen the weighted node-packing model, and \citet{Kroon1997} show that the problem is NP-complete as soon as each train has three routing possibilities. \citet{Caprara2011} study a general formulation with a quadratic objective, and \citet{Lusby2011} develop a resource-based set-packing formulation and a branch-and-price algorithm for junctions. \citet{Lusby2011survey} survey track allocation models at the strategic, tactical, and operational levels. These studies mainly decide whether, and how well, a given timetable can be accommodated in a station. In integrated planning, a timetable change also alters the occupation intervals, and changing a stop to a pass alters the admissible station routing alternatives.

\subsection{Integrated timetabling and station track and route assignment}
Many studies integrate train scheduling and station track and route assignment. Decomposition is a common strategy: \citet{Lamorgese2015} couple line and station dispatching through feasibility cuts, \citet{Lamorgese2016} develop a Benders-like reformulation connecting the two infrastructure levels, and \citet{Lamorgese2017} extend the approach to cyclic and non-cyclic timetabling. For microscopic timetable planning, \citet{Leutwiler2022} propose an LBBD with logical feasibility cuts generated by a SAT-based procedure, extended with set covering heuristics by \citet{Leutwiler2023}. \citet{Besinovic2016} coordinate microscopic running and headway information with macroscopic schedules, and \citet{Wang2023} jointly optimize train scheduling and routing on a coupled multi-resolution space--time network with Lagrangian relaxation and branch-and-bound.

Direct formulations of the integrated problem also exist. \citet{Petering2016} jointly formulate cyclic timetabling and platforming, \citet{Burggraeve2017} build a routing plan and a cyclic timetable for a complex station area, and \citet{Caimi2011} propose a resource-tree conflict graph model for conflict-free routing and scheduling in stations. \citet{Xu2017} and \citet{Zhou2016} integrate train routing and timetabling on lines with station layouts, using MILP and train-based Lagrangian decomposition, respectively. Closest to our decision scope, \citet{Xu2021} consider stop-skipping, platform choice, and passenger flows on a one-way line with a time--space network flow formulation and a Lagrangian heuristic. Maintenance adds further coupling: \citet{VanAken2017} adjust timetables under infrastructure possessions, \citet{ZhangY2019a} integrate microscopic timetabling and track maintenance scheduling, \citet{ZhangQ2020} re-optimize timetables and platform schedules under planned maintenance, \citet{ZhangQ2022} integrate timetabling, platforming, and maintenance scheduling in a mesoscopic model with dynamically updated train time windows, and \citet{Ji2024} adjust stop patterns, train orders, and track assignment at key stations under uncertain maintenance plans.

In real-time traffic management, the alternative graph of \citet{Mascis2002} is widely used to represent train scheduling with blocking resources \citep{DAriano2007}; \citet{Pellegrini2014} formulate routing and scheduling in complex junctions as a MILP, and \citet{Sama2017} re-route trains with variable neighbourhood search. These approaches model the station infrastructure at a fine level, but stop decisions and service frequency requirements are not part of the problem.

\subsection{Large neighborhood search and MILP-based matheuristics}
Large neighborhood search (LNS) improves a solution by repeatedly destroying part of it and repairing the released part \citep{Shaw1998}. ALNS adds adaptive selection among several destroy and repair heuristics \citep{Ropke2006,Pisinger2007}; \citet{Mara2022} survey its applications and hybridizations. In railway scheduling, ALNS has been used for stop planning and timetabling \citep{Dong2020} and for demand-driven timetabling \citep{Yin2024}, typically with heuristic removal and insertion of trains. When the repair step is itself a MILP over the released variables, the method belongs to the family of MIP-based neighborhood searches, such as local branching \citep{Fischetti2003} and relaxation induced neighborhoods \citep{Danna2005}. Conflicts can also organize the search: \citet{Reisch2021} resolve large conflict graphs of freight train slots by local search. A difficulty that is specific to our problem is that a valid neighborhood must contain not only the decisions in conflict, but also the decisions that compensate for the resulting changes in service counts.

\subsection{Research gap}
\Cref{tab:lit} summarizes the decisions and methods of representative studies. Integrated studies of stop planning and timetabling usually treat stations macroscopically, whereas integrated studies of timetabling and platforming usually take the stopping pattern as given. \citet{Xu2021} and \citet{Ji2024} come closest to our decision scope; to our knowledge, however, no study combines station service frequency requirements evaluated on time windows with track occupation and route headways in station throats. Moreover, existing ALNS approaches for railway timetabling remove and reinsert trains heuristically, while decomposition approaches return feasibility information from stations to a master problem. This paper addresses the gap by (i) representing stops, timetables, and station routings on one event-activity network with disjunctive and conditional activities, and (ii) designing neighborhoods that release station conflicts together with their service frequency compensation and repairing them exactly by a restricted MILP.

\begin{table}[tb]
\centering\small
\caption{Decision scope and solution methods of representative studies (\checkmark: considered).}\label{tab:lit}
\setlength{\tabcolsep}{3.5pt}
\begin{tabular}{@{}>{\raggedright\arraybackslash}p{32mm}cccc>{\raggedright\arraybackslash}p{58mm}@{}}\toprule
Study & Stops & Freq. & Tracks/routes & Maint. & Model; solution method\\\midrule
\citet{Yang2016} & \checkmark & \checkmark & & & MILP; commercial solver\\
\citet{Dong2020} & \checkmark & \checkmark & & & Nonlinear model; extended ALNS\\
\citet{Yao2023} & \checkmark & & & & Periodic time--space network; ADMM\\
\citet{Lusby2011} & & & \checkmark & & Set packing; branch-and-price\\
\citet{Lamorgese2016} & & & \checkmark & & MILP; Benders-like decomposition\\
\citet{Burggraeve2017} & & & \checkmark & & ILP/MILP; sequential routing and timetabling\\
\citet{Xu2021} & \checkmark & & platforms & & Network flow; Lagrangian heuristic\\
\citet{ZhangQ2020} & & & \checkmark & \checkmark & Space--time network; Lagrangian relaxation, ADMM\\
\citet{Leutwiler2022} & & & micro. & & MILP; logic-based Benders decomposition\\
\citet{Wang2023} & & & \checkmark & & Space--time network; Lagrangian relaxation, branch-and-bound\\
\citet{Ji2024} & \checkmark & & tracks & \checkmark & Two-stage stochastic program; dual decomposition\\
This paper & \checkmark & \checkmark & \checkmark & \checkmark & Event-activity network MILP; CD-ALNS with MILP repair\\\bottomrule
\end{tabular}
\end{table}
